\section*{Capítulo \ref{ch:b1:precipitation} --- Precipitación y solubilidad}

\begin{solution}{exo:b1:precipitation:1}
\ce{BaSO4(s) <=> Ba^2+ + SO4^2-}, $K_s = [\ce{Ba^2+}][\ce{SO4^2-}] =
10^{-9.97}$. \ce{CaF2(s) <=> Ca^2+ + 2F-}, $K_s = [\ce{Ca^2+}][\ce{F-}]^2 =
10^{-9.84}$. \ce{Ag2CrO4(s) <=> 2Ag+ + CrO4^2-}, $K_s =
[\ce{Ag+}]^2[\ce{CrO4^2-}] = 10^{-11.95}$. \ce{Fe(OH)3(s) <=> Fe^3+ + 3OH-},
$K_s = [\ce{Fe^3+}][\ce{OH-}]^3 = 10^{-38.55}$. Para $K_s = \num{2.0e-13}$,
$\mathrm{p}K_s = 12.70$.
\end{solution}

\begin{solution}{exo:b1:precipitation:2}
$s = \sqrt{K_s} = 10^{-12.27/2} = \qty{7.3e-7}{mol/L}$, es decir,
$\num{7.3e-7} \times 187.8 = \qty{1.4e-4}{g/L} = \qty{0.14}{mg/L}$.
\end{solution}

\begin{solution}{exo:b1:precipitation:3}
Tras la mezcla (volumen duplicado), $[\ce{Ba^2+}] = \qty{1.0e-3}{mol/L}$ y
$[\ce{SO4^2-}] = \qty{1.0e-4}{mol/L}$: $Q_s = \num{1.0e-7} > K_s =
\num{1.1e-10}$: precipita el sulfato de bario. En el segundo caso,
$[\ce{SO4^2-}] = \num{1.0e-7}$ y $Q_s = \num{1.0e-10} < K_s$: no hay
precipitado, aunque por muy poco.
\end{solution}

\begin{solution}{exo:b1:precipitation:4}
En agua, $s = 10^{-9.97/2} = \qty{1.0e-5}{mol/L}$. En sulfato de
\qty{0.010}{mol/L} (\cref{prop:b1:precipitation:common-ion}),
$s = 10^{-9.97}/0.010 =
\qty{1.1e-8}{mol/L}$, unas mil veces menos (factor 970).
\end{solution}

\begin{solution}{exo:b1:precipitation:5}
En agua, $s = (10^{-9.84}/4)^{1/3} = \qty{3.3e-4}{mol/L}$, es decir,
\qty{26}{mg/L}. En cloruro de calcio de \qty{0.10}{mol/L},
$[\ce{Ca^2+}] \approx 0.10$ y $[\ce{F-}] = 2s$, de modo que
$s = \frac12\sqrt{10^{-9.84}/0.10} = \qty{1.9e-5}{mol/L}$. La
aproximación $s \ll 0.10$ se cumple con cuatro órdenes de magnitud de
margen.
\end{solution}

\begin{solution}{exo:b1:precipitation:6}
Comienzo: $\mathrm{pH} = 14.00 - \frac12(11.25 + \log 0.050) = 14.00 - 4.97 =
9.03$. Al \qty{99}{\percent}, $[\ce{Mg^2+}] = \num{5.0e-4}$ y
$\mathrm{pH} = 14.00 - \frac12(11.25 - 3.30) = 10.03$. En un eje de pH,
\ce{Mg(OH)2} existe por encima de 9.03; por debajo, el magnesio está
íntegramente como \ce{Mg^2+} a \qty{0.050}{mol/L}.
\end{solution}

\begin{solution}{exo:b1:precipitation:7}
El bromuro de plata aparece primero, cuando $[\ce{Ag+}] = 10^{-12.27 + 2} =
10^{-10.27}$, y el cloruro de plata, en $10^{-7.75}$. Cuando aparece el
cloruro de plata, $[\ce{Br-}] = 10^{-12.27}/10^{-7.75} = 10^{-4.52} =
\qty{3.0e-5}{mol/L}$, el \qty{0.30}{\percent} del bromuro. Con el criterio
habitual del \qty{0.1}{\percent}, la separación no es aceptable: los dos
$\mathrm{p}K_s$ solo difieren en 2.5.
\end{solution}

\begin{solution}{exo:b1:precipitation:8}
El cromato de plata se forma cuando $[\ce{Ag+}]^2 \times \num{5.0e-3} = 10^{-11.95}$,
$[\ce{Ag+}] = \qty{1.5e-5}{mol/L}$. Entonces, $[\ce{Cl-}] = 10^{-9.75}/\num{1.5e-5}
= \qty{1.2e-5}{mol/L}$, una fracción \num{2.4e-4} (\qty{0.024}{\percent})
del cloruro: el color rojo aparece cuando prácticamente todo el cloruro
ha precipitado.
\end{solution}

\begin{solution}{exo:b1:precipitation:9}
$[\ce{CO3^2-}] = [\ce{HCO3-}]K_{a2}/[\ce{H3O+}] = [\ce{HCO3-}]\,10^{\mathrm{pH}
- \mathrm{p}K_{a2}}$, luego $K_s = [\ce{Ca^2+}][\ce{HCO3-}]\,10^{\mathrm{pH} -
\mathrm{p}K_{a2}}$, que es la relación pedida. Con $[\ce{Ca^2+}] =
[\ce{HCO3-}] = s$: $s^2 = 10^{-8.30 + 10.33 - 8.30} = 10^{-6.27}$,
$s = \qty{7.3e-4}{mol/L}$ (\qty{73}{mg/L} de \ce{CaCO3}), frente a
$10^{-4.15} = \qty{7.1e-5}{mol/L}$ si el ion carbonato no fuera una base.
\end{solution}

\begin{solution}{exo:b1:precipitation:10}
1.~Sumando \ce{CO2(g) <=> CO2(aq)}: $K' = K K_H = 10^{-4.34 - 1.47} =
10^{-5.81}$.
2.~$[\ce{Ca^2+}] = s$ y $[\ce{HCO3-}] = 2s$, de modo que $4s^3 = K'p$. A
\qty{0.010}{bar}: $s = (10^{-5.81} \times 0.010/4)^{1/3} =
\qty{1.6e-3}{mol/L}$, \qty{157}{mg/L}. Con el aire exterior:
$s = \qty{5.5e-4}{mol/L}$, \qty{55}{mg/L}. (Comprobación: con el aire del
suelo, $[\ce{CO2(aq)}] = 10^{-3.47}$ y $\mathrm{pH} = 6.37 + \log(3.14\times
10^{-3}/3.39 \times 10^{-4}) = 7.34$: el hidrogenocarbonato predomina
efectivamente.)
3.~El agua saturada de calcita a la elevada presión de dióxido de carbono
del suelo llega al techo de la cueva, donde la presión es menor: pierde
dióxido de carbono, $Q > K'$, y la calcita precipita donde cuelga la
gota, anillo tras anillo.
% ledger: air:CO2, dfg:CO2_g, dfg:CO2_ao (log KH computed in figdata/bachelor-1/aqueous-constants.py)
\end{solution}

\begin{solution}{exo:b1:precipitation:11}
1.~Comienzo en $\mathrm{pH} = 14.00 - \frac12(16.38 - 2.00) = 6.81$.
Redisolución completa cuando $[\ce{Zn(OH)4^2-}] = 10^{-2} = 10^{-1.93}[\ce{OH-}]^2$,
$[\ce{OH-}] = 10^{-0.035}$, $\mathrm{pH} = 13.97$. El hidróxido existe
entre pH 6.81 y 13.97.
2.~$\mathrm{pH} = 14.00 - 14.45/4 = 10.39$, $s_{\min} = 2 \times
10^{-(16.38 + 1.93)/2} = \qty{1.4e-9}{mol/L}$ (en este modelo de dos
especies).
3.~Recta de pendiente $-2$ desde $(6.81, -2)$ hasta el mínimo, hacia
$(10.39, -8.9)$, y después pendiente $+2$ hasta $(13.97, -2)$; la curva
real queda redondeada por las demás especies hidroxo
(figura de la \cref{sec:b1:precipitation:amphoteric}).
\end{solution}

\begin{solution}{exo:b1:precipitation:12}
1.~$\ce{Al(OH)3(s) + OH- <=> Al(OH)4-}$, $K = 10^{-1.23}$; todo el
aluminio está disuelto cuando $10^{-3} \leqslant K[\ce{OH-}]$, es decir,
$[\ce{OH-}] \geqslant 10^{-1.77}$, $\mathrm{pH} \geqslant 12.23$.
2.~$[\ce{Fe^3+}] = 10^{-38.55}/10^{-3 \times 1.77} = 10^{-33.2}$: el
hierro permanece íntegramente como \ce{Fe(OH)3}; la filtración separa el
hierro (sólido) del aluminio (filtrado).
3.~El hidróxido de magnesio no es anfótero: en exceso de hidróxido sigue
sólido, junto con el hidróxido de hierro(III), y nada se separa.
\end{solution}

\begin{solution}{pb:b1:precipitation:1}
\textbf{1.} \ce{Fe(OH)3(s) <=> Fe^3+ + 3OH-}, $K_s = [\ce{Fe^3+}][\ce{OH-}]^3$;
\ce{Zn(OH)2(s) <=> Zn^2+ + 2OH-}, $K_s = [\ce{Zn^2+}][\ce{OH-}]^2$;
\ce{Mg(OH)2(s) <=> Mg^2+ + 2OH-}, análogamente.
\textbf{2.} Las constantes no tienen la misma forma (potencia 3 o 2 de
$[\ce{OH-}]$) y los tres iones están a concentraciones distintas; solo
pueden compararse los pH de comienzo.
\textbf{3.} Al comienzo, $c[\ce{OH-}]^n = K_s$, de modo que
$n\log[\ce{OH-}] =
-\mathrm{p}K_s - \log c$ y $\mathrm{pH} = \mathrm{p}K_e +
\log[\ce{OH-}] = \mathrm{p}K_e - \frac1n(\mathrm{p}K_s + \log c)$.
\textbf{4.} $14.00 - \frac13(38.55 - 3.00) = 2.15$.
\textbf{5.} $14.00 - \frac12(16.38 - 2.30) = 6.96$.
\textbf{6.} $14.00 - \frac12(11.25 - 1.70) = 9.22$.
\textbf{7.} El hierro(III) (2.15), después el cinc (6.96) y después el
magnesio (9.22). A pH 2.0, $Q_s = 10^{-3} \times 10^{-36} = 10^{-39} < 10^{-38.55}$
para el hierro, y los demás están aún más lejos de precipitar: nada está
en estado sólido.
\textbf{8.} \qty{1.0e-6}{mol/L}.
\textbf{9.} $14.00 - \frac13(38.55 - 6.00) = 3.15$.
\textbf{10.} $\log[\ce{Fe^3+}] = -38.55 + 3(14.00 - \mathrm{pH}) = 3.45 -
3\,\mathrm{pH}$: pendiente $-3$.
\textbf{11.} Hasta el comienzo de la precipitación del hidróxido de cinc,
pH 6.96.
\textbf{12.} De 3.15 a 6.96.
\textbf{13.} \ce{Fe^3+ + OH- <=> FeOH^2+}, $\beta_1 =
[\ce{FeOH^2+}]/([\ce{Fe^3+}][\ce{OH-}]) = 10^{11.82}$.
\textbf{14.} Las dos son iguales cuando $[\ce{OH-}] = 10^{-11.82}$, a pH
2.18: \ce{Fe^3+} predomina por debajo, y \ce{FeOH^2+}, por encima.
\textbf{15.} $10^{11.82 + 3.15 - 14.00} = 10^{0.97} = 9.3$: a pH 3.15, el
hierro disuelto es 10.3 veces el \ce{Fe^3+} libre, \qty{1.0e-5}{mol/L}:
solo se elimina el \qty{99.0}{\percent}. La primera estimación no era
segura.
\textbf{16.} $[\ce{Fe}]_{\mathrm{disuelto}} = 10^{3.45 - 3\,\mathrm{pH}}\,
(1 + 10^{\mathrm{pH} - 2.18})$.
\textbf{17.} $\log[\ce{FeOH^2+}] = 11.82 + \log[\ce{Fe^3+}] + \mathrm{pH} -
14.00 = 1.27 - 2\,\mathrm{pH}$: pendiente $-2$.
\textbf{18.} $1.27 - 2\,\mathrm{pH} \approx -6.00$ da 3.64; incluyendo el
término de \ce{Fe^3+}, $10^{3.45 - 10.92}(1 + 10^{1.46}) =
\num{1.0e-6}$ a $\mathrm{pH} = 3.64$.
\textbf{19.} $[\ce{Zn^2+}] = \num{5.0e-5}$: $14.00 - \frac12(16.38 - 4.30) =
7.96$. El hidróxido de magnesio empieza en 9.22: no.
\textbf{20.} \ce{Zn(OH)2(s) + 2OH- <=> Zn(OH)4^2-}, $K = \beta_4 K_s =
10^{14.45 - 16.38} = 10^{-1.93}$.
\textbf{21.} $[\ce{OH-}]^2 = \num{5.0e-3}/10^{-1.93}$, $\mathrm{pH} =
13.81$. El operador no debe pasarse con la base: por encima de pH 9.2 el
magnesio precipita con el cinc y, mucho más arriba, el cinc vuelve a la
disolución.
\textbf{22.} $\num{1.0e-3} \times 1000 \times 0.999 \times 106.8 =
\qty{107}{g}$ por metro cúbico.
\textbf{23.} $\num{5.0e-3} \times 1000 \times 0.99 \times 99.4 =
\qty{492}{g}$ por metro cúbico.
\textbf{24.} Si se sube de golpe por encima de 6.96, los dos hidróxidos
precipitan en un mismo lodo: el cinc se pierde en el residuo de hierro, y
el residuo queda contaminado con cinc. Detenerse entre los umbrales
elimina solo el hierro.
\textbf{25.} \textbf{De pH 3.64 a pH 6.96}: por debajo de 3.64 sigue
disuelto más del \qty{0.1}{\percent} del hierro, sobre todo como
\ce{FeOH^2+}; por encima de 6.96 se forma hidróxido de cinc. Despreciar
\ce{FeOH^2+} habría situado el límite inferior en 3.15.
\end{solution}
